% from main.tex
%%%%%%%%%%  scan idx 166  |  folio 153  %%%%%%%%%%
\textbf{Corollaire 7.7.} \emph{Il existe, pour} \(E\in\Ob\Dpx(S)_{\underline{\mathrm{parf}}}\) et \(F\in\Ob\Dpx^+(S)\), \emph{un isomorphisme fonctoriel}
\[
  E^\vee\otL F\xrightarrow{\sim}\mathbb{R}\underline{\Hom}(E,F),
\]
\emph{avec la notation de} (7.3).

\emph{Preuve.} Appliquer (7.6), \(F\) et \(G\) étant remplacés respectivement par \(A\) et \(F\), et tenir compte du fait (7.1) que \(E^\vee\in\Ob\Dpx(S)_{\underline{\mathrm{parf}}}\).

\textbf{Corollaire 7.8.} \emph{Soient} \(E,F\in\Ob\Dpx(S)_{\underline{\mathrm{parf}}}\). \emph{Il existe une flèche naturelle}
\[
  \mathbb{R}\underline{\Hom}(E,F)\otL\mathbb{R}\underline{\Hom}(F,E)\longrightarrow A,
\]
\emph{qui est une «~dualité parfaite~», i.e.\ identifie modulo (7.4) chacun des complexes} \(\mathbb{R}\underline{\Hom}(E,F)\), \(\mathbb{R}\underline{\Hom}(F,E)\) \emph{au dual de l'autre (au sens de (7.3))}.

\emph{Preuve.} D'après (7.5), on a des accouplements naturels \(E^\vee\otL E\to A\), \(F^\vee\otL F\to A\), d'où, par produit tensoriel, un accouplement naturel
\[
  E^\vee\otL F\otL F^\vee\otL E\longrightarrow A,
\]
qui, compte tenu de (7.7), donne l'accouplement annoncé sur les \(\mathbb{R}\underline{\Hom}\). Montrons qu'il s'agit d'une dualité parfaite. On a les isomorphismes canoniques
\[
\begin{aligned}
\mathbb{R}\underline{\Hom}(F^\vee\otL E,A)&\simeq\mathbb{R}\underline{\Hom}(F^\vee,E^\vee) &&(7.4)\\
&\simeq F^{\vee\vee}\otL E^\vee &&(7.6)\\
&\simeq F\otL E^\vee &&(7.2).
\end{aligned}
\]
Donc, moyennant une vérification de compatibilité, on a gagné. On pourrait aussi se localiser et dévisser pour se ramener à \(E=F=A\).


